\chapter{绪论}

\section{研究背景与意义}

\subsection{研究背景}

随着信息技术的快速发展和网络规模的持续扩大，网络系统面临的攻击手段日益复杂。传统的异常检测方法主要依赖人工编制的规则库，面对未知攻击时检出率较低，且规则维护成本较高。因此，研究能够快速发现异常流量、具备良好泛化能力的检测方法具有重要的理论意义和实际应用价值。

网络异常检测的本质是从大量正常行为中识别出偏离预期的模式\cite{liu1997}。深度学习方法能够自动学习数据中的高层次特征表示，无需人工设计复杂的特征规则，已成为该领域的重要研究方向。

\subsection{研究意义}

本文的研究意义主要体现在两个方面：在理论上，通过引入深度自编码器与注意力机制，探索高维流量数据的紧凑表示方法；在工程上，设计并实现一个完整的检测系统，为校园网络安全运维提供可落地的工具支撑。

\section{国内外研究现状}

国内外学者围绕网络异常检测开展了大量研究。早期工作多基于统计分析方法和机器学习方法，近年来深度学习方法逐渐成为主流。文献\cite{yuan2001}采用神经网络对流量特征进行分类，取得了较好的检测效果。总体来看，现有研究在检测精度方面已有较好基础，但在模型轻量化与实时性方面仍有改进空间。

\section{论文主要工作与组织结构}

本文的主要工作包括：第一，构建流量特征数据集并完成预处理；第二，设计基于深度自编码器的异常检测模型；第三，实现系统各功能模块并进行测试验证。

全文共分为四章。第 1 章为绪论，介绍研究背景、意义与论文结构；第 2 章介绍相关理论与技术；第 3 章给出系统设计与实现细节并对实验结果进行分析；最后是结论、参考文献、致谢与附录。
